% ================= APÉNDICE =================
\appendix
\chapter{Los 25 verdadero/falso clásicos}
Todos salieron (con esta u otra redacción) en parciales entre 2007 y 2022. Tapá la columna de la derecha. $A,B,C$ son cuadradas $n\times n$ salvo que se diga otra cosa.

{\small
\renewcommand{\arraystretch}{1.3}
\begin{tabularx}{\linewidth}{@{}r X c@{}}
\toprule
& \textbf{Afirmación} & \\
\midrule
1 & $AB=AC$ y $A\neq O\Rightarrow B=C$. & \textcolor{rojo}{\textbf{F}}\\
2 & $(A+B)^2=A^2+2AB+B^2$. & \textcolor{rojo}{\textbf{F}}\\
3 & $AB=O\Rightarrow A=O$ o $B=O$. & \textcolor{rojo}{\textbf{F}}\\
4 & $A$, $B$ invertibles $\Rightarrow A+B$ invertible. & \textcolor{rojo}{\textbf{F}}\\
5 & $\lambda\neq0$, $A,B$ invertibles $\Rightarrow(\lambda AB)^{-1}=\frac1\lambda B^{-1}A^{-1}$. & \textcolor{verde}{\textbf{V}}\\
6 & $\det(2A)=2\det A$. & \textcolor{rojo}{\textbf{F}}\\
7 & $\det(A+B)=\det A+\det B$. & \textcolor{rojo}{\textbf{F}}\\
8 & $A^2=I\Rightarrow A$ es invertible. & \textcolor{verde}{\textbf{V}}\\
9 & $\det A=0\Rightarrow A$ tiene una fila o una columna de ceros. & \textcolor{rojo}{\textbf{F}}\\
10 & $AA^t$ es simétrica y $\det(AA^t)\ge0$ ($A$ real). & \textcolor{verde}{\textbf{V}}\\
11 & Un sistema con más incógnitas que ecuaciones siempre tiene solución. & \textcolor{rojo}{\textbf{F}}\\
12 & Un sistema con más incógnitas que ecuaciones, si tiene solución, no es única. & \textcolor{verde}{\textbf{V}}\\
13 & El rango de $A$ es la cantidad de filas no nulas de $A$. & \textcolor{rojo}{\textbf{F}}\\
14 & Toda matriz elemental tiene determinante 1. & \textcolor{rojo}{\textbf{F}}\\
15 & $AB$ invertible $\iff A$ y $B$ invertibles. & \textcolor{verde}{\textbf{V}}\\
16 & Si $X_1,X_2$ resuelven $AX=b$, entonces $\frac12(X_1+X_2)$ también. & \textcolor{verde}{\textbf{V}}\\
17 & Antisimétrica con $n$ impar $\Rightarrow\det A=0$. & \textcolor{verde}{\textbf{V}}\\
18 & $\rg(AB)=n\Rightarrow\rg A=\rg B=n$. & \textcolor{verde}{\textbf{V}}\\
19 & Toda invertible es producto de elementales, y toda elemental es invertible. & \textcolor{verde}{\textbf{V}}\\
20 & $AX=0$ es SCD $\Rightarrow A$ invertible y $A^{-1}b$ resuelve $AX=b$ para todo $b$. & \textcolor{verde}{\textbf{V}}\\
21 & $A$ de $5\times7$ $\Rightarrow$ existe $b$ con $AX=b$ compatible determinado. & \textcolor{rojo}{\textbf{F}}\\
22 & Si $A$ no es cuadrada, $\rg(A)\neq\rg(A^t)$. & \textcolor{rojo}{\textbf{F}}\\
23 & Si $E$ es una forma escalonada de $A$, entonces $\det E=\det A$. & \textcolor{rojo}{\textbf{F}}\\
24 & $A=PBP^{-1}$ con $P$ invertible $\Rightarrow$ ($A$ invertible $\iff B$ invertible). & \textcolor{verde}{\textbf{V}}\\
25 & $\tr(AB)=\tr A\cdot\tr B$. & \textcolor{rojo}{\textbf{F}}\\
\bottomrule
\end{tabularx}}

\begin{memo}{CONTRAEJEMPLOS DE BOLSILLO}
$N=\begin{pmatrix}0&1\\0&0\end{pmatrix}$: $N\neq O$ pero $N^2=O$. Rompe 1 ($NN=NO$), 3 y la cancelación.\\
$A=I$, $B=-I$: rompe 4 y 7 ($\det O=0$, pero $1+(-1)^n$ puede ser 2).\\
$\begin{pmatrix}1&1\\1&1\end{pmatrix}$: sin ceros, determinante 0, rango 1 (rompe 9 y 13).\\
$E=\mathrm{diag}(1,5)$: elemental con determinante 5 (rompe 14).\\
23: escalerizar cambia el determinante (un intercambio cambia el signo, escalar una fila lo multiplica).
\end{memo}

\chapter{Formulario en una página}
{\small
\begin{memo}{COMPLEJOS}
$z\conj z=|z|^2$ \quad $\frac1z=\frac{\conj z}{|z|^2}$ \quad $|zw|=|z||w|$ \quad $e^{i\theta}=\cos\theta+i\sin\theta$\\
$re^{i\theta}\cdot se^{i\varphi}=rse^{i(\theta+\varphi)}$ \quad $(re^{i\theta})^n=r^ne^{in\theta}$ \quad $z^n=Re^{i\varphi}$: $z_k=\sqrt[n]R\,e^{i(\varphi+2k\pi)/n}$, $k=0,\dots,n-1$
\end{memo}
\begin{memo}{SISTEMAS}
Rouché--Frobenius: compatible $\iff\rg A=\rg(A|b)$; \ SCD si además $\rg A=n$; SCI con $n-\rg A$ parámetros.\\
Sol. general = una particular + las del homogéneo. \ $m<n$: nunca SCD. \ Homogéneo: siempre compatible.
\end{memo}
\begin{memo}{MATRICES}
$(AB)^t=B^tA^t$ \quad $(AB)^{-1}=B^{-1}A^{-1}$ \quad $(A^t)^{-1}=(A^{-1})^t$ \quad $(\lambda A)^{-1}=\frac1\lambda A^{-1}$\\
$\tr(AB)=\tr(BA)$ \quad $\begin{pmatrix}a&b\\c&d\end{pmatrix}^{-1}=\frac1{ad-bc}\begin{pmatrix}d&-b\\-c&a\end{pmatrix}$ \quad $(A|I)\to(I|A^{-1})$\\
$(A+B)^2=A^2+AB+BA+B^2$ \quad $\rg A=\rg A^t\le\min(m,n)$ \quad $\rg(PA)=\rg A$ si $P$ invertible
\end{memo}
\begin{memo}{DETERMINANTES}
$F_i\leftrightarrow F_j$: $\times(-1)$ \quad $F_i\to\lambda F_i$: $\times\lambda$ \quad $F_i\to F_i+\lambda F_j$: igual \quad $\det A^t=\det A$\\
$\det(AB)=\det A\det B$ \quad $\det(\lambda A)=\lambda^n\det A$ \quad $\det A^{-1}=\frac1{\det A}$ \quad triangular: producto de la diagonal\\
Invertible $\iff\det\neq0\iff\rg=n\iff AX=0$ solo trivial\\
Antisim. $n$ impar: 0 \quad nilpotente: 0 \quad $A^2=A$: 0 o 1 \quad $A^2=I$: $\pm1$ \quad $\det\begin{pmatrix}A&B\\O&C\end{pmatrix}=\det A\det C$
\end{memo}
\begin{trampa}{ANTES DE MARCAR}
¿Sustituiste el valor crítico y terminaste de escalerizar? \ ¿$\det(\lambda A)$ con $\lambda^n$? \ ¿Factorizaste respetando izquierda y derecha? \ ¿Chequeaste la inversa multiplicando una fila por una columna? \ ¿Dibujaste el complejo para ver el cuadrante?
\end{trampa}}
